\chapter{La parete diventa un ingresso}
\label{ch:wallinput}

Finora il solutore ha disegnato da sé la propria parete di ugello:
la forma usciva dal calcolo, e l'unico modo di cambiarla era
cambiare le poche manopole che guidano la costruzione. Un
ottimizzatore che debba esplorare forme libere ha bisogno del flusso
di dati rovesciato --- gli si consegna una parete qualunque e se ne
ottiene il flusso su quella parete. Questo capitolo esegue quella
inversione usando gli stessi mattoni già verificati, confronta il
nuovo motore con quello vecchio sull'unica parete che entrambi
sanno trattare, e certifica la derivata della spinta rispetto a una
deformazione arbitraria della parete. È quella derivata la bussola
di ogni ottimizzazione successiva di questo documento.

\section{La domanda}

La marcia del Mattone~1 \emph{disegna} la propria parete: il
contorno è calcolato come una \emph{linea di corrente} --- una linea
ovunque tangente alla velocità locale, il percorso seguito da una
particella di gas --- scelta in modo che al suo interno passi
esattamente la portata di gola. In altre parole la parete è
un'uscita. Il progetto a forma libera richiede l'opposto: si
consegna una parete \emph{qualsiasi} e se ne ottiene il flusso. Le
stesse celle già verificate sanno calcolare il flusso su una parete
assegnata?

\begin{figure}[htbp]
\centering
\includegraphics[width=0.86\textwidth]{sch_inversion.png}
\caption{L'inversione strutturale, disegnata in coordinate di ugello
($x$ [m] lungo l'asse, $y$ [m] raggio). A sinistra: il generatore
decide da sé la propria parete seguendo una linea di corrente di
portata, quindi la forma è un risultato. A destra: l'ottimizzatore
consegna una parete arbitraria e le celle di parete inversa ne
calcolano il flusso, quindi la forma è un dato. Si guardino le
frecce: ciò che cambia è soltanto la direzione in cui vengono
determinate la posizione di parete e lo stato del flusso --- tutto
ciò che sta a monte, e ogni cella interna al flusso, resta
invariato. È per questo che il nuovo motore eredita la
certificazione del vecchio.}
\label{fig:inversion}
\end{figure}

\section{Come funziona, per esteso}

Il motore lascia invariato tutto ciò che sta a monte della parete
--- stesso avvio in gola, stesso ventaglio di espansione iniziale ---
e cambia soltanto il modo in cui ciascuna colonna di punti calcolati
termina. Nel generatore il punto di parete è \emph{prodotto}: la
marcia integra la linea di corrente di portata e la parete è
ovunque essa vada. In modalità inversa la \emph{posizione} del punto
di parete è \emph{data} (deve trovarsi sul contorno assegnato in
quella stazione, con la pendenza locale del contorno fornita come
dato), quindi la cella di parete lavora all'indietro, esattamente
come descritto nel Capitolo~\ref{ch:moc}: trova il \emph{piede}
della \emph{caratteristica} in arrivo --- il punto sulla colonna
precedente da cui è partita la linea che porta informazione al punto
di parete --- con una ricerca per corde, scandendo i segmenti
memorizzati di quella colonna alla ricerca di quello i cui dati
racchiudono la condizione caratteristica. Chiude poi la relazione di
compatibilità lungo la corda dal piede al punto di parete, con la
pendenza assegnata della parete che vi fornisce la direzione del
flusso~\cite{zucrow}. Si costruisce una colonna di marcia per ogni
stazione di parete assegnata.

Due conseguenze danno forma al resto del motore. Primo, una parete
arbitraria non produce un'uscita uniforme e assiale, quindi nulla a
valle può assumerlo: la spinta è sempre valutata per la via di
parete~\eqref{eq:route2}, che richiede solo le pressioni di parete e
il flusso in gola. Secondo, l'intera costruzione resta
differenziabile --- la ricerca del piede, la chiusura sulla corda,
la quadratura della spinta --- quindi la differenziazione automatica
in modo inverso fornisce $\mathrm{d}J/\mathrm{d}({\rm parete})$, la
sensibilità della spinta a una variazione di parete a forma libera,
al costo di circa una marcia aggiuntiva.

Quel gradiente è certificato dall'\emph{identità dell'inviluppo},
che è il teorema fondamentale del calcolo usato come test: la media
di una derivata su un intervallo deve eguagliare la pendenza della
secante su quell'intervallo,
$\overline{\mathrm{d}J/\mathrm{d}a} = [J(a_2)-J(a_1)]/(a_2-a_1)$. Il
test deforma dunque la parete con una protuberanza liscia e
localizzata di ampiezza $a$, spazza $a$ su un intervallo, adatta una
retta ai campioni del gradiente AD e richiede che la pendenza media
adattata coincida con una secante ampia di $J$ stesso entro una
banda derivata. Perché non l'ovvio confronto puntuale fra AD e una
differenza finita a passo piccolo? Perché le proprietà del gas sono
lette da tabelle interpolate linearmente a tratti, quindi la
derivata esatta del codice porta un microscopico dente di sega alla
scala delle celle di tabella, e una differenza a passo piccolo si
mette a cavallo di quelle cuspidi: quel test è qui mal posto. Il
risultato e la sua diagnosi sono nell'Appendice~\ref{app:ripple}; il
certificato d'inviluppo ne è la soluzione e la
Figura~\ref{fig:ripple} il riassunto in un'immagine.

\begin{figure}[htbp]
\centering
\includegraphics[width=0.72\textwidth]{fig_grad_ripple.pdf}
\caption{Perché il gradiente è certificato per inviluppo anziché
puntualmente. Asse orizzontale: ampiezza $a$ [m] della protuberanza
liscia aggiunta alla parete, da rientrante a sporgente. Asse
verticale: la sensibilità della spinta $\mathrm{d}J/\mathrm{d}a$
[N/m] restituita dalla differenziazione automatica, cioè newton di
spinta guadagnati per metro di protuberanza. Si guardino le due cose
sovrapposte: un andamento rettilineo e liscio (la retta; curvatura
adattata $J'' = -7.3\times10^{7}$~N/m$^2$), che è la fisica, e una
minuscola dispersione attorno a esso (barre rosse a $4\sigma$;
$\sigma = 44$~N/m, lo 0.010\% del gradiente), che è il motivo di
cuspidi delle tabelle di proprietà lineari a tratti. Entrambi sono
proprietà reali dell'obiettivo discreto --- l'AD riporta la derivata
esatta del codice, dente di sega compreso --- ma solo l'andamento è
fisico, ed è per questo che il certificato media sull'increspatura
invece di campionarla. Storia completa
nell'Appendice~\ref{app:ripple}.}
\label{fig:ripple}
\end{figure}

\section{Che cosa è stato costruito}

\carrier{validation/a1\_wall\_march.py}: il motore a parete
assegnata (\code{wall\_march}), l'estrattore di stazioni
(\code{stations\_from\_wall}), la spinta per via di parete
(\code{J\_wall}), il flusso di massa per colonna, il certificato di
gradiente per identità dell'inviluppo e i rigettatori (una parete i
cui dati di pendenza siano corrotti deve far fallire la marcia).

\section{Risultati di riferimento}

\textbf{Stadio gemello, 5/5 PASS.} Alimentato con il contorno del
Mattone~1 come parete assegnata, il nuovo motore riproduce la spinta
del Mattone~1 attraverso un calcolo completamente diverso ---
$J_{\rm wall} = 1.25534\times10^{8}$~N contro
$J_{\rm exit} = 1.25521\times10^{8}$~N, una differenza di
$1.37\times10^{4}$ dentro una banda derivata di $1.58\times10^{5}$.
La portata attraverso l'ultima colonna calcolata eguaglia la portata
di gola a meno di 12.3~kg/s (banda 152), una verifica di
conservazione indipendente che il confronto fra gemelli non
contiene; il numero di Mach minimo lungo la parete è 1.20, quindi il
flusso resta supersonico ovunque e il metodo resta applicabile;
17{,}221 celle sono certificate da Newton, cioè ogni risoluzione
locale è convergita dentro la propria banda derivata di residuo; e
la parete con pendenza corrotta fa fallire la marcia --- rigettata.
\textbf{Stadio gradiente, 2/2 PASS} per identità dell'inviluppo:
$|{\rm AD} - {\rm secante}| = 29$ contro una banda derivata di 109,
e la sensibilità è essa stessa molto maggiore di quella banda,
quindi la parete conta in modo dimostrabile.
